% ============================================================
%  SOS/A FORMELSAMMLUNG — Teil 3–5
%  Zustandsraum & ÜF | Nichtlineare Systeme | Linearisierung & Lyapunov
% ============================================================

\bereich[cKap3]{Teil 3: Zustandsraumdarstellung und Übertragungsfunktion}

\begin{formelreihe}
  \formelblock{Zustandsraum (allgemein)}{\begin{aligned}
    \dot{\mathbf{x}}(t) &= \mathbf{A}\,\mathbf{x}(t)+\mathbf{B}\,u(t) \\
    y(t) &= \mathbf{C}\,\mathbf{x}(t)+D\,u(t)
  \end{aligned}} &
  \infoblock{$\mathbf{A}$ ($n\!\times\!n$): Systemmatrix\\[0.05em]
    $\mathbf{B}$ ($n\!\times\!1$): Eingangsmatrix\\[0.05em]
    $\mathbf{C}$ ($1\!\times\!n$): Ausgangsmatrix\\[0.05em]
    $D$: Durchgriff (Skalar)} &
  \beispielblock{RC-Glied}{%
    \dot{x}=-\tfrac{1}{RC}x+\tfrac{1}{RC}u\\[0.05em]
    y=x\\[0.05em]
    A=-\tfrac{1}{RC},\;B=\tfrac{1}{RC}} \\[0.3em]

  \formelblock{Übertragungsfunktion aus ZR}{\begin{aligned}
    G(s) = \mathbf{C}(s\mathbf{I}-\mathbf{A})^{-1}\mathbf{B}+D
  \end{aligned}} &
  \parbox[t]{\linewidth}{%
    \formelblock{Charakterist. Polynom}{%
      N(s)=\det(s\mathbf{I}-\mathbf{A})}\\[0.25em]%
    \infoblock{EW von $\mathbf{A}$ = Pole von $G(s)$\\
      $\Rightarrow$ Routh auf $N(s)$ anwendbar\\[0.04em]
      $n\!=\!2$: $N\!=\!s^2-\operatorname{tr}(\mathbf{A})\,s+\det(\mathbf{A})$}} &
  \beispielblock{3×3 Sarrus-Regel}{%
    \mathbf{A}\!=\!\begin{bsmallmatrix}1&2&0\\3&0&1\\0&4&2\end{bsmallmatrix}\\[0.07em]
    {\scriptstyle\begin{array}{@{}c@{\,}c@{\,}c@{\,}|@{\,}c@{\,}c@{}}
      (s\!-\!1)&-2&0&(s\!-\!1)&-2\\[-0.02em]
      -3&s&-1&-3&s\\[-0.02em]
      0&-4&(s\!-\!2)&0&-4
    \end{array}}\\[0.07em]
    =s^3\!-\!3s^2\!-\!8s\!+\!16} \\[0.3em]

  \formelblock{Faltung (Zeit)}{%
    y(t)=\int_0^t h(\tau)\,u(t-\tau)\,\mathrm{d}\tau} &
  \formelblock{Multiplikation (Frequenz)}{%
    Y(s)=G(s)\cdot U(s)\\[0.08em]
    G(s)=\dfrac{Y(s)}{U(s)}} &
  \beispielblock{Impulsantwort}{%
    h(t)=\mathcal{L}^{-1}\{G(s)\}\\[0.05em]
    \text{PT1: }h(t)=\tfrac{K}{T}e^{-t/T}} \\
\end{formelreihe}
\fallstrick{$G(s)$ setzt \textbf{LTI} voraus: $Y\!=\!G\,U$ gilt nur bei linear + zeitinvariant. Nichtlin.\ System $\Rightarrow$ \textbf{kein} globales $G(s)$ (erst um Ruhelage linearisieren, dann nur lokal).}
\bereichende

% ============================================================

\bereich[cKap4]{Teil 4: Nichtlineare Systeme — Ruhelagen \& Stabilitätsanalyse}

\begin{formelreihe}
  \formelblock{Nichtlin. System (allg.)}{\begin{aligned}
    \dot{\mathbf{x}}(t) &= \mathbf{f}(\mathbf{x}(t),u(t),t)\\
    y(t) &= h(\mathbf{x}(t),u(t),t)
  \end{aligned}} &
  \parbox[t]{\linewidth}{%
    \formelblock{Autonomes System}{%
      \dot{\mathbf{x}}(t)=\mathbf{f}(\mathbf{x}(t))\\
      \text{(zeitinvariant, }u=0\text{)}}\\[0.25em]%
    \infoblock{NL-Besonderheiten: mehrere isolierte\\
      Ruhelagen, Hysterese, Chaos möglich}} &
  \beispielblock{Pendel}{\begin{aligned}
    \dot{x}_1 &= x_2\\
    \dot{x}_2 &= -\tfrac{g}{\ell}\sin x_1
  \end{aligned}} \\[0.3em]

  \formelblock{Ruhelage $\mathbf{x}_R$}{\begin{aligned}
    &\mathbf{f}(\mathbf{x}_R,\,\mathbf{0})=\mathbf{0}\\[0.06em]
    &\text{1. }u(t)=0\\
    &\text{2. }\dot{x}_i\stackrel{!}{=}0\\
    &\text{3. }x_i\to x_{i,R}\\
    &\text{4. Löse }\mathbf{f}(\mathbf{x}_R,\mathbf{0})=\mathbf{0}
  \end{aligned}} &
  \infoblock{Autonomes System: kein\\
    Eingang ($u\equiv 0$).\\[0.05em]
    NL: mehrere isolierte\\
    Ruhelagen möglich.} &
  \beispielblock{Pendel-RL}{%
    \dot{x}_1\!=\!x_2\stackrel{!}{=}0
    \Rightarrow x_{2,R}\!=\!0\\[0.04em]
    \dot{x}_2\!=\!-\tfrac{g}{\ell}\sin x_1\stackrel{!}{=}0\\[0.02em]
    \Rightarrow\sin x_{1,R}\!=\!0\\[0.02em]
    \Rightarrow x_{1,R}\in\{0,\pi\}\\[0.05em]
    \text{RL: }(0,0),\,(\pi,0)} \\[0.3em]

  \formelblock{Betriebspunkt $\mathbf{x}_B$}{%
    u_B=\mathrm{const},\quad
    \mathbf{f}(\mathbf{x}_B,\,u_B)=\mathbf{0}} &
  \infoblock{$\mathbf{x}_B$ $=$ RL des\\
    linearisierten $\Delta$-Systems\\
    (Teil\,5:
    $\Delta\mathbf{x}=\mathbf{0}$
    $\Leftrightarrow\mathbf{x}=\mathbf{x}_B$)} &
  \beispielblock{BP: $\dot{x}=u-x^2$}{%
    u_B=1:\;x_B=\pm1\\[0.04em]
    (2\text{ Betriebspunkte})} \\[0.3em]

  \formelblock{Vektornormen}{\begin{aligned}
    \|\mathbf{x}\|_1 &= \textstyle\sum_i|x_i|\\
    \|\mathbf{x}\|_2 &= \sqrt{\textstyle\sum_i x_i^2}\\
    \|\mathbf{x}\|_\infty &= \max_i|x_i|
  \end{aligned}} &
  \infoblock{Norm-Eigenschaften:\\[0.05em]
    $\|\mathbf{x}\|=0\Leftrightarrow\mathbf{x}=\mathbf{0}$\\[0.03em]
    $\|c\,\mathbf{x}\|=|c|\,\|\mathbf{x}\|$\\[0.03em]
    Dreiecksungl.:
    $\|\mathbf{x}+\mathbf{y}\|\leq\|\mathbf{x}\|+\|\mathbf{y}\|$} &
  \beispielblock{2-Norm}{%
    \mathbf{x}=\begin{bmatrix}3\\4\end{bmatrix}\\[0.1em]
    \|\mathbf{x}\|_2=\sqrt{9+16}=5} \\[0.3em]

  \parbox[t]{\linewidth}{%
    \formelblock{Stabil nach Lyapunov (iSvL)}{%
      \forall\varepsilon>0\;\exists\,\delta>0:\\
      \|\mathbf{x}(0)-\mathbf{x}_R\|<\delta\\
      \Rightarrow\|\mathbf{x}(t)-\mathbf{x}_R\|<\varepsilon\;\forall t\geq0}\\[0.2em]%
    {\scriptsize\scshape\color{gray!65!black} Wann iSvL? (Kriterien)}\\[0.03em]%
    \scriptsize
    \textbf{Anschaulich:} kleine Auslenk.\ $\Rightarrow$ Traj.\ bleibt \emph{nahe} RL, kehrt \textbf{nicht} zwingend zurück; Grenzfall asym.\,stab.$\leftrightarrow$instab.\\[0.08em]
    \textbf{Lyapunov:} $\exists\,V$ pos.\,def.\ mit $\dot V\leq0\;\forall\mathbf{x}$ (Satz\,1).\\[0.05em]
    \textbf{EW (LZI/linear.):} alle $\operatorname{Re}(\lambda_i)\leq0$ \emph{und} EW auf $j\omega$-Achse \textbf{einfach}.\\[0.05em]
    \textbf{Typisch:} rein imag.\ EW $\pm j\omega$ (ungedämpfte Schwingung, konst.\ Amplitude) $\Rightarrow$ iSvL, aber \emph{nicht} asym.} &
  \parbox[t]{\linewidth}{%
    \formelblock{Asymptotisch stabil}{%
      \text{iSvL-stabil PLUS:}\\
      \lim_{t\to\infty}\|\mathbf{x}(t)-\mathbf{x}_R\|=0}\\[0.25em]%
    \formelblock{Global asymptotisch stabil}{%
      \text{Asym. stabil für alle}\\
      \mathbf{x}(0)\in\mathbb{R}^n}\\[0.15em]%
    \infoblock{\textbf{Absolut stabil:} glob.\,asym.\,stabil\\
      mit $\mathbf{x}_R=\mathbf{0}$ (RL im Nullpunkt)}} &
  \beispielblock{Stabilitäten}{%
    \text{Pendel m. Reib.: asym. stab.}\\[0.08em]
    \text{PT1: glob. asym. stab.}\\[0.08em]
    \text{Ungedämpft. Pendel /}\\
    \text{Oszillator, EW }\pm j\omega:\\
    \text{nur iSvL (grenzstabil)}\\[0.08em]
    \text{EW mit Re}>0:\text{ instabil}} \\
\end{formelreihe}
\fallstrick{\textbf{Hat das System einen Eingang $u$?}\quad
\textbullet\ \textbf{JA} (auch $G(s)$, $G(j\omega)$): nur \textbf{BIBO} (Pole/EW alle $\operatorname{Re}<0$?). Lyapunov-Ruhelage $=$ ``nicht anwendbar, da nicht autonom''.\quad
\textbullet\ \textbf{NEIN} ($\dot{\mathbf{x}}\!=\!\mathbf{A}\mathbf{x}$, kein $u$): nur \textbf{Lyapunov-Ruhelage} (EW alle $\operatorname{Re}<0$). BIBO $=$ ``nicht anwendbar''.}
\bereichende

% ============================================================

\bereich[cKap4]{Rezept: Ruhelagen \& Betriebspunkt bestimmen}

\begin{rezept}
\rs{\textbf{Autonom?} Kein Eingang in den Gln.\ $\Rightarrow$ autonom. Ruhelage: $u\equiv0$;\; Betriebspunkt: $u\equiv u_B=$ const.}
\rs{Alle $\dot x_i\stackrel{!}{=}0$ setzen (stationär).}
\rs{Gleichungssystem $\mathbf{f}(\mathbf{x}_R,\mathbf{0})=\mathbf{0}$ bzw.\ $\mathbf{f}(\mathbf{x}_B,u_B)=\mathbf{0}$ nach $x_{i}$ lösen.\;
  \textbf{Mehrere} Lösungen $=$ mehrere isolierte Ruhelagen $\Rightarrow$ \emph{alle} angeben.}
\end{rezept}

\noindent\footnotesize\textbf{Betriebspunkt aus Soll-Ausgang $y_B$} (typ.\ Klausurform) $=$ Gleichungssystem aus $y_B{=}h(\mathbf{x}_B,u_B)$ \& allen $\dot x_i{=}0$; schrittweise lösen:\quad
\textbullet\ $y_B\to x_{jB}$\quad
\textbullet\ $\dot x_2{=}0\to$ weitere $x_{iB}$\quad
\textbullet\ $\dot x_1{=}0\to u_B$.\par\vspace{0.1em}

\miniexample{$\dot x_1=-x_1+2+u(x_1^2-2)$,\; $\dot x_2=2x_1-x_2(3-x_1)$,\; $y=2x_2$;\; Soll $y_B=4$}{%
$y_B=2x_{2B}=4\Rightarrow x_{2B}=2$.\quad
$\dot x_2=0:\;2x_{1B}-2(3-x_{1B})=0\Rightarrow x_{1B}=1{,}5$.\quad
$\dot x_1=0:\;-1{,}5+2+u_B(1{,}5^2-2)=0\Rightarrow u_B=-2$.}

\fallstrick{$\sin,\cos$ im \textbf{Bogenmaß (rad)} auswerten!\quad
\textbullet\ \emph{alle} isolierten Ruhelagen angeben.\quad
\textbullet\ Betriebspunkt $\mathbf{x}_B$ wird im $\Delta$-System zur Ruhelage $\Delta\mathbf{x}=\mathbf{0}$.}

\bereichende

\bereich[cKap5]{Transformation in Betriebspunkt}

\begin{formelreihe}
  \formelblock{Vorgehen (3 Schritte)}{\begin{aligned}
    &\text{1. }\dot{x}_i\stackrel{!}{=}0,\;u=u_B
      \;\Rightarrow x_{iB}\\[0.08em]
    &\text{2. Substitution:}\\
    &\quad x_i=\Delta x_i+x_{iB},\\
    &\quad u=\Delta u+u_B\\
    &\quad\text{in Originalgl.\ einsetzen}\\
    &\quad\Rightarrow\Delta\dot{x}_i\text{-Gleichungen}\\[0.08em]
    &\text{3. }x_{iB},\,u_B\text{ einsetzen;}\\
    &\quad\text{Terme ohne }\Delta\text{ kürzen}\\
    &\quad(\,f(\mathbf{x}_B,u_B)=\mathbf{0}\,)
  \end{aligned}} &
  \infoblock{Terme ohne $\Delta$ fallen weg\\
    (im BP: $f(\mathbf{x}_B,u_B)=\mathbf{0}$)\\[0.1em]
    Ergebnis: $\Delta$-System,\\
    RL $\Delta\mathbf{x}=\mathbf{0}$\\
    $(\Leftrightarrow\mathbf{x}=\mathbf{x}_B)$\\[0.05em]
    \textup{= Trafo in die RL}\\[0.07em]
    \textbf{Nachweis:} $\Delta x_i\!=\!0\;\forall i$:\\
    r.\,S.\,$=0\Rightarrow\Delta\mathbf{x}_R\!=\!\mathbf{0}$\\[0.08em]
    \textbf{Exakt} — NL-Terme\\
    in $\Delta x_i$ bleiben!} &
  \beispielblock{Bsp.\ $u_B=2$}{%
    \dot{x}_1=-x_1+2u\\
    \dot{x}_2=x_1^2-x_2-4\\[0.05em]
    x_{1B}=4,\;x_{2B}=12\\[0.08em]
    \Delta\dot{x}_1=-\Delta x_1+2\Delta u\\
    \Delta\dot{x}_2=8\Delta x_1\\
    \quad+(\Delta x_1)^2-\Delta x_2} \\
\end{formelreihe}
\bereichende

% ============================================================

\bereich[cKap5]{Linearisierung via Jacobi-Matrix}

\begin{formelreihe}
  \formelblock{Vorgehen}{\begin{aligned}
    &\text{1. }\dot{x}_i=f_i(x_1,\ldots,x_n,u)\\
    &\quad\text{aufschreiben}\\[0.07em]
    &\text{2. Partiell ableiten:}\\
    &\quad A_{B,ij}=\tfrac{\partial f_i}{\partial x_j}\big|_{BP}\\[0.02em]
    &\quad b_{B,i}=\tfrac{\partial f_i}{\partial u}\big|_{BP}\\[0.07em]
    &\text{3. }x_{jB},\,u_B\text{ einsetzen}\\[0.07em]
    &\text{4. }\mathbf{A}_B,\,\mathbf{b}_B\text{ aufstellen}
  \end{aligned}} &
  \parbox[t]{\linewidth}{%
    \formelblock{Matrixform $(n\!=\!2)$}{\begin{aligned}
      \mathbf{A}_B &=
        \begin{bsmallmatrix}
          A_{B,11}&A_{B,12}\\[0.04em]
          A_{B,21}&A_{B,22}
        \end{bsmallmatrix},\;
      \mathbf{b}_B =
        \begin{bsmallmatrix}b_{B,1}\\[0.04em]b_{B,2}\end{bsmallmatrix}
    \end{aligned}}\\[0.25em]
    \infoblock{$A_{B,ij}=\tfrac{\partial f_i}{\partial x_j}\big|_{BP}$\\
      (im BP ausgewertet).\\[0.06em]
      \textbf{Einsetzen nie vergessen!}}} &
  \beispielblock{Pendel, $u_B\!=\!0$}{%
    {\scriptstyle f_2=(u\cos x_1-g\sin x_1-dx_2)/\ell}\\[0.06em]
    \text{stehend }(x_{1B}=\pi)\text{:}\\[0.02em]
    \mathbf{A}_B\!=\!\begin{bsmallmatrix}0&1\\+g/\ell&-d/\ell\end{bsmallmatrix}\!,\;
    \mathbf{b}_B\!=\!\begin{bsmallmatrix}0\\-1/\ell\end{bsmallmatrix}\\[0.1em]
    \text{hängend }(x_{1B}=0)\text{:}\\[0.02em]
    \mathbf{A}_B\!=\!\begin{bsmallmatrix}0&1\\-g/\ell&-d/\ell\end{bsmallmatrix}\!,\;
    \mathbf{b}_B\!=\!\begin{bsmallmatrix}0\\+1/\ell\end{bsmallmatrix}} \\[0.3em]

  \formelblock{Taylor-Linearisierung}{\begin{aligned}
    \Delta\dot{x}_i &\approx
    \textstyle\sum_j\frac{\partial f_i}{\partial x_j}\Delta x_j
    +\frac{\partial f_i}{\partial u}\Delta u\\
    &\text{(ausgewertet bei }(\mathbf{x}_B,u_B)\text{)}
  \end{aligned}} &
  \formelblock{Linearisiertes System}{\begin{aligned}
    \Delta\dot{\mathbf{x}} &= \mathbf{A}_B\,\Delta\mathbf{x}+\mathbf{b}_B\,\Delta u\\
    \Delta y &= \mathbf{C}_B\,\Delta\mathbf{x}+d_B\,\Delta u
  \end{aligned}} &
  \beispielblock{1. Ordnung}{%
    \dot{x}=u-x^3+2,\;u_B=2\\[0.04em]
    x_B=0\\[0.04em]
    A_B=0,\;b_B=1\\[0.04em]
    \Delta\dot{x}=\Delta u} \\[0.3em]

  \parbox[t]{\linewidth}{%
    \formelblock{Jacobi-Matrizen}{\begin{aligned}
      \mathbf{A}_B &=
        \tfrac{\partial\mathbf{f}}{\partial\mathbf{x}}
        \big|_{(\mathbf{x}_B,u_B)},\;
      \mathbf{b}_B =
        \tfrac{\partial\mathbf{f}}{\partial u}
        \big|_{(\mathbf{x}_B,u_B)}\\[0.08em]
      \mathbf{C}_B &=
        \tfrac{\partial h}{\partial\mathbf{x}}
        \big|_{(\mathbf{x}_B,u_B)},\;
      d_B =
        \tfrac{\partial h}{\partial u}
        \big|_{(\mathbf{x}_B,u_B)}
    \end{aligned}}\\[0.25em]
    \formelblock{Delta-Koordinaten}{\begin{aligned}
      \mathbf{x} &= \Delta\mathbf{x}+\mathbf{x}_B,\quad
      u = \Delta u+u_B\\
      \Delta\mathbf{x} &= \mathbf{0}
        \Leftrightarrow \mathbf{x}=\mathbf{x}_B
    \end{aligned}}} &
  \parbox[t]{\linewidth}{%
    {\scriptsize\scshape\color{gray!65!black} Stabilität via Linearisierung}\\[0.03em]%
    1.\ $\mathbf{A}_B$ aufstellen (Jacobi)\\[0.02em]
    2.\ EW: $\det(s\mathbf{I}-\mathbf{A}_B)=0$ lösen $\Rightarrow s_\infty$\\[0.02em]
    3.\ RT aller $s_\infty$ prüfen:\\[0.01em]
    {\scriptsize
    \quad alle neg.\,RT $\Rightarrow$ as.\,stab.\,(lok.)\\[0.01em]
    \quad pos.\,RT vorh. $\Rightarrow$ instab.\\[0.03em]
    \quad lokal (nahe BP); LZI: global\\[0.01em]
    \quad neg.\,RT: kehrt zurück\\
    \quad pos.\,RT: läuft weg}} &
  \beispielblock{Trick $2\!\times\!2$, $\operatorname{tr}=0$}{%
    \mathbf{A}_B=\begin{bsmallmatrix}0&1\\\alpha&0\end{bsmallmatrix}\!:\\[0.06em]
    s^2\!-\!\alpha=0\Rightarrow s_\infty=\pm\sqrt{\alpha}\\[0.05em]
    \alpha>0\Rightarrow\text{instab.}\\[0.02em]
    \alpha<0\Rightarrow\text{grenzstab.\ (rein imag.)}\\[0.06em]
    \text{Pendel stehend: }\alpha=\tfrac{g}{\ell}>0\\
    \Rightarrow\text{instab.}} \\
\end{formelreihe}
\bereichende

% ============================================================

\bereich[cKap5]{Rezept: Linearisierung $\to$ Kleinsignal-ZR $(\mathbf{A},\mathbf{b},\mathbf{c}^T,d)$ + BIBO}

\begin{rezept}
\rs{Betriebspunkt $(\mathbf{x}_B,u_B)$ bestimmen (s.\ oben).}
\rs{Jacobi-Matrizen bilden:\;
  $\mathbf{A}=\tfrac{\partial\mathbf{f}}{\partial\mathbf{x}}$,\;
  $\mathbf{b}=\tfrac{\partial\mathbf{f}}{\partial u}$,\;
  $\mathbf{c}^T=\tfrac{\partial h}{\partial\mathbf{x}}$,\;
  $d=\tfrac{\partial h}{\partial u}$\;
  (Ausgangsgl.\ $y=h$ nicht vergessen!).}
\rs{Zahlenwerte $(\mathbf{x}_B,u_B)$ \textbf{einsetzen} $\Rightarrow$ konstante $\mathbf{A},\mathbf{b},\mathbf{c}^T,d$.}
\rs{Kleinsignal-ZR:\;
  $\Delta\dot{\mathbf{x}}=\mathbf{A}\Delta\mathbf{x}+\mathbf{b}\Delta u$,\;
  $\Delta y=\mathbf{c}^T\Delta\mathbf{x}+d\Delta u$.}
\rs{\textbf{BIBO:} EW von $\mathbf{A}$ aus $\det(s\mathbf{I}-\mathbf{A})=0$ (=Pole). Alle $\operatorname{Re}<0\Rightarrow$ stabil; ein EW mit $\operatorname{Re}\geq0\Rightarrow$ instabil. (Routh statt EW-Rechnung möglich.)}
\end{rezept}

\miniexample{$\mathbf{f}$ wie oben, BP $x_{1B}=x_{2B}=u_B=1$,\; $h=2x_2$}{%
$\tfrac{\partial f_1}{\partial x_1}=-1+2ux_1=1$,\; $\tfrac{\partial f_1}{\partial x_2}=0$,\; $\tfrac{\partial f_1}{\partial u}=x_1^2-2=-1$;\quad
$\tfrac{\partial f_2}{\partial x_1}=2+x_2=3$,\; $\tfrac{\partial f_2}{\partial x_2}=-(3-x_1)=-2$,\; $\tfrac{\partial f_2}{\partial u}=0$.\quad
$\Rightarrow\mathbf{A}=\begin{bsmallmatrix}1&0\\3&-2\end{bsmallmatrix}$,\;
$\mathbf{b}=\begin{bsmallmatrix}-1\\0\end{bsmallmatrix}$,\;
$\mathbf{c}^T=[\,0\;\;2\,]$,\; $d=0$.\quad
EW: $s=1,\,-2$ $\Rightarrow$ ein EW $+1$ $\Rightarrow$ \textbf{instabil}.}

\bereichende

\bereich[cKap5]{Teil 5: Linearisierung im Betriebspunkt \& Direkte Methode von Lyapunov}

\begin{formelreihe}
  \parbox[t]{\linewidth}{%
    {\scriptsize\scshape\color{gray!65!black} Direkte Methode von Lyapunov}\\[0.03em]%
    1.\ $V(\mathbf{x})$ wählen: pos.\,def., $V(\mathbf{0})=0$\\[0.02em]
    2.\ $\dot{V}=\bigl(\tfrac{\partial V}{\partial\mathbf{x}}\bigr)^{\!T}\!\mathbf{f}(\mathbf{x})$ berechnen\\[0.02em]
    3.\ Definitheit von $\dot{V}$ prüfen:\\[0.01em]
    {\scriptsize
    \quad $\dot{V}\leq0\Rightarrow$ glob.\,stab.\ (Satz\,1)\\[0.01em]
    \quad $\dot{V}<0\Rightarrow$ glob.\,asym.\,stab.\ (Satz\,2)}} &
  \infoblock{Typische $V$-Ansätze:\\[0.03em]
    $V=\mathbf{x}^T\mathbf{Q}\mathbf{x}$,\ $\mathbf{Q}$ pos.\,def.\\[0.02em]
    $V=\sum_i\lambda_i x_i^2$\ ($\lambda_i>0$)} &
  \beispielblock{$\dot{x}=-x^3$}{%
    V=\tfrac{1}{2}x^2>0\\[0.04em]
    \dot{V}=x(-x^3)=-x^4\\[0.04em]
    \dot{V}<0\;\forall x\!\neq\!0\\
    \Rightarrow\text{gl.\,asym.\,stab.}} \\[0.3em]

  \formelblock{Definitheit (Funktion $V$)}{\begin{aligned}
    \text{pos. def.}&:\;V(\mathbf{x})>0\;\forall\mathbf{x}\neq\mathbf{0}\\
    \text{pos. semid.}&:\;V(\mathbf{x})\geq0\\
    \text{neg. def.}&:\;V(\mathbf{x})<0\;\forall\mathbf{x}\neq\mathbf{0}\\
    \text{neg. semid.}&:\;V(\mathbf{x})\leq0\\
    \text{indefinit}&:\;\text{sonst (Vorz.-Wechsel)}\\
    &V(\mathbf{0})=0\;\text{in allen Fällen}
  \end{aligned}} &
  \parbox[t]{\linewidth}{%
    {\scriptsize\scshape\color{gray!65!black} Definitheit (Matrix $\mathbf{Q}$)}\\[0.03em]%
    \scriptsize
    \textbf{Voraussetzung:} $\mathbf{Q}$ symmetrisch.\\[0.02em]
    \quad$\mathbf{Q}$ \textbf{nicht sym.}\ $\Rightarrow$ \textbf{indefinit!} (s.\ $\mathbf{Q}_5$)\\[0.05em]
    \textbf{EW-Kriterium} ($\mathbf{Q}$ sym., alle EW reell):\\[0.02em]
    \quad alle EW$>0$: pos.\,def.;\quad alle EW$<0$: neg.\,def.\\[0.02em]
    \quad alle EW$\geq0$: pos.\,semid.;\; alle EW$\leq0$: neg.\,semid.\\[0.02em]
    \quad Vorzeichenwechsel der EW $\Rightarrow$ indefinit\\[0.05em]
    \textbf{Hauptunterdet.-Krit.}\ $D_k=\det\mathbf{Q}_k$ (führend):\\[0.02em]
    \quad pos.\,def.: alle $D_k>0$ \ ($q_{11}>0,\,D_2>0,\dots$)\\[0.02em]
    \quad neg.\,def.: $(-1)^k D_k>0$ \ ($q_{11}<0,\,D_2>0,\dots$)\\[0.02em]
    \qquad äquiv.\ zu $-\mathbf{Q}$ pos.\,def.\ (s.\ $\mathbf{Q}_2$)\\[0.02em]
    \quad semidef.: alle $D_k\geq0$ bzw.\ $\leq0$, Randfall $\det\mathbf{Q}=0$\\[0.02em]
    \quad indefinit: sonst (z.\,B.\ $\det\mathbf{Q}<0$ bei $2\!\times\!2$)\\[0.05em]
    $V=\mathbf{x}^T\!\mathbf{Q}\mathbf{x}$ erbt Definitheit von $\mathbf{Q}$.} &
  \beispielblock{Definitheit prüfen}{%
    \mathbf{Q}_1=\begin{bsmallmatrix}2&-1\\-1&2\end{bsmallmatrix}:\;
    q_{11}{=}2{>}0,\,\det{=}3{>}0\\
    \Rightarrow\text{pos.\,def.}\\[0.1em]
    \mathbf{Q}_2=\begin{bsmallmatrix}-1&-1\\-1&-2\end{bsmallmatrix}:\;
    -\mathbf{Q}_2{=}\begin{bsmallmatrix}1&1\\1&2\end{bsmallmatrix}\\
    1{>}0,\,\det{=}1{>}0\Rightarrow-\mathbf{Q}_2\text{ p.\,d.}\\
    \Rightarrow\mathbf{Q}_2\;\text{neg.\,def.}\\[0.1em]
    \mathbf{Q}_3=\begin{bsmallmatrix}-1&1\\1&-1\end{bsmallmatrix}:\;
    q_{11}{<}0,\,\det{=}0\\
    \Rightarrow\text{neg.\,semidef.}\\[0.1em]
    \mathbf{Q}_4=\begin{bsmallmatrix}1&2\\2&1\end{bsmallmatrix}:\;
    \det{=}1{-}4{=}{-}3{<}0\\
    \Rightarrow\text{indefinit}\\[0.1em]
    \mathbf{Q}_5=\begin{bsmallmatrix}3&1\\2&2\end{bsmallmatrix}\text{ \textbf{unsym.}}
    \Rightarrow\text{indefinit}} \\[0.3em]

  \formelblock{Satz 1: global stabil}{\begin{aligned}
    &V(\mathbf{x})\text{ pos. def., }V(\mathbf{0})=0\\
    &\dot{V}(\mathbf{x})\leq0\;\forall\mathbf{x}\neq\mathbf{0}\\
    &\Rightarrow\text{Ruhelage global stabil}
  \end{aligned}} &
  \formelblock{Satz 2: global asympt. stabil}{\begin{aligned}
    &V(\mathbf{x})\text{ pos. def., }V(\mathbf{0})=0\\
    &\dot{V}(\mathbf{x})<0\;\forall\mathbf{x}\neq\mathbf{0}\\
    &\Rightarrow\text{global asym. stabil}
  \end{aligned}} &
  \beispielblock{Satz 2 Bsp}{%
    \dot{x}=-x:\\[0.05em]
    V=\tfrac{1}{2}x^2>0\\[0.05em]
    \dot{V}=x(-x)=-x^2<0\\
    \Rightarrow\text{gl. asym. stab.}} \\[0.3em]

  \formelblock{$\dot{V}$ berechnen}{\begin{aligned}
    \dot{V}(\mathbf{x})
    &= \sum_i\frac{\partial V}{\partial x_i}\dot{x}_i\\[0.05em]
    &= \left(\frac{\partial V}{\partial\mathbf{x}}\right)^{\!T}
       \mathbf{f}(\mathbf{x})
  \end{aligned}} &
  \infoblock{\textbf{Einzugsgebiet:} alle
    $\mathbf{x}(0)$ mit
    $\mathbf{x}(t)\to\mathbf{x}_R$ ($t\to\infty$).\\[0.1em]
    \textbf{Nachweis:} Inneres einer
    Äqui-$V$-Linie $V(\mathbf{x}){=}c$
    ganz in $\{\dot{V}{<}0\}$.} &
  \beispielblock{Einzugsgebiet}{%
    \dot{x}_1=-x_1+2x_1^3 x_2\\
    \dot{x}_2=-x_2\\[0.05em]
    V=\lambda_1 x_1^2+\lambda_2 x_2^2\\
    \text{Ellipse als EZG}} \\[0.3em]

  \parbox[t]{\linewidth}{%
    {\scriptsize\scshape\color{gray!65!black}%
      Stabilität LZI — Zusammenfassung}\\[0.03em]%
    \scriptsize\itshape
    BIBO $\Leftrightarrow$ alle Pole neg. Realteil\\[0.03em]
    iSvL glob. stabil $\Leftrightarrow$ keine EW pos. RT,\\
    \phantom{iSvL glob. stabil}\,keine mehrfachen EW auf $j\omega$\\[0.03em]
    iSvL glob. asym. stab. $\Leftrightarrow$ alle EW neg. RT} &
  \parbox[t]{\linewidth}{%
    \scriptsize\itshape
    Glob. stab. (Lyap.): $V$ pos.def.\ $+$ $\dot{V}\!\leq\!0$\\[0.03em]
    Glob. asym. (Lyap.): $V$ pos.def.\ $+$ $\dot{V}\!<\!0$\\[0.06em]
    Via Linearisierung \textbf{(lokal!):}\\[0.03em]
    stab. RL: keine EW pos. RT,\\
    \phantom{stab. RL:}\,keine mehrfachen EW auf $j\omega$\\[0.03em]
    asym. stab. RL: alle EW neg. RT} &
  \beispielblock{Merkregel}{%
    \dot{V}\leq0\text{: stabil}\\
    \dot{V}<0\text{: asym. stab.}\\[0.05em]
    \text{(beide global)}\\[0.05em]
    \text{Lin.: nur lokal!}} \\
\end{formelreihe}
\bereichende

% ============================================================

\bereich[cKap5]{Rezept: Definitheit mit Parametern $(a,b)$ — Hauptunterdeterminanten}

\begin{rezept}
\rs{$\mathbf{Q}$ \textbf{symmetrisch}? Sonst indefinit. ($P=\mathbf{x}^T\mathbf{Q}\mathbf{x}$ erbt Definitheit von $\mathbf{Q}$; nicht-diagonale Symmetrie kann Parameter festlegen.)}
\rs{Führende Hauptunterdet.\ $D_1=q_{11},\;D_2=\det\mathbf{Q}_2,\;\dots,\;D_n=\det\mathbf{Q}$.}
\rs{\textbf{pos.\ def.:} alle $D_k>0$.\quad
  \textbf{neg.\ def.:} $(-1)^k D_k>0$ (Vorz.\ $-,+,-,\dots$), äquiv.\ $-\mathbf{Q}$ pos.\ def.}
\rs{Ungleichungen nach $a,b$ auflösen.}
\end{rezept}

\miniexample{$\mathbf{Q}=\protect\begin{bsmallmatrix}a&1\\1&b\protect\end{bsmallmatrix}$ (2$\times$2) \;und\; $\mathbf{Q}=\protect\begin{bsmallmatrix}a&2&0\\2&3&2\\b&2&4\protect\end{bsmallmatrix}$ (3$\times$3)}{%
\textbf{2$\times$2:} pos.\ def.\ $\Leftrightarrow a>0$ \emph{und} $D_2=ab-1>0$ ($b>\tfrac1a$);\;
neg.\ def.\ $\Leftrightarrow a<0$ \emph{und} $ab-1>0$ ($b<\tfrac1a$).\\[0.15em]
\textbf{3$\times$3:} Symmetrie $q_{13}=q_{31}\Rightarrow b=0$. Dann $D_1=a$, $D_2=3a-4$, $D_3=8a-16$;\;
pos.\ def.\ $\Leftrightarrow a>2$ (und $b=0$).}

\bereichende

\bereich[cKap5]{Rezept: Direkte Methode von Lyapunov \& Einzugsgebiet}

\begin{rezept}
\rs{$V(\mathbf{x})$ pos.\ def.\ wählen, $V(\mathbf{0})=0$\;
  (Ansatz $V=\mathbf{x}^T\mathbf{Q}\mathbf{x}$ mit $\mathbf{Q}$ pos.\ def., oder $V=\textstyle\sum_i\lambda_i x_i^2$, $\lambda_i>0$).}
\rs{$\dot V=\bigl(\tfrac{\partial V}{\partial\mathbf{x}}\bigr)^{\!T}\mathbf{f}(\mathbf{x})=\textstyle\sum_i\tfrac{\partial V}{\partial x_i}\dot x_i$ berechnen, $\dot x_i$ einsetzen.}
\rs{Definitheit von $\dot V$ prüfen:\;
  $\dot V\leq0\;\forall\mathbf{x}\Rightarrow$ global \emph{stabil} (Satz\,1);\;
  $\dot V<0\;\forall\mathbf{x}\neq\mathbf{0}\Rightarrow$ global \emph{asym.}\ stabil (Satz\,2).}
\end{rezept}

\noindent\footnotesize\textbf{Einzugsgebiet} (lokal, wenn $\dot V$ \emph{nicht} überall neg.\ ist):\par\vspace{0.05em}
\begin{rezept}
\rs{$\dot V$ nach Quadrat-Termen sortieren $\Rightarrow$ Gebiet $\{\dot V<0\}$ als Bedingung an $x_i$.}
\rs{Größte Niveaumenge $\{V(\mathbf{x})<c\}$, die \textbf{ganz} in $\{\dot V<0\}$ liegt $=$ nachgewiesenes Einzugsgebiet.}
\rs{In $x_1$-$x_2$-Ebene skizzieren (Kreis bei $V=\sum x_i^2$, Ellipse bei $V=\sum\lambda_i x_i^2$).}
\end{rezept}

\miniexample{$\dot x_1=-x_1$,\; $\dot x_2=-x_2+5x_1x_2$,\; $V=\tfrac12 x_1^2+\tfrac12 x_2^2$}{%
$\dot V=x_1(-x_1)+x_2(-x_2+5x_1x_2)=-x_1^2-x_2^2(1-5x_1)$.\quad
$\dot V<0$ falls $1-5x_1>0$, d.\,h.\ $x_1<\tfrac15$.\quad
Größter Kreis um $\mathbf{0}$ darin: $r=\tfrac15=0{,}2$ $\Rightarrow$ EZG $=$ Kreis mit $r=0{,}2$.}

\fallstrick{$\dot V\leq0$ (semidef.)\ liefert nur \emph{stabil}, \textbf{nicht} asymptotisch.\quad
\textbullet\ Lyapunov ist nur \textbf{hinreichend}: kein Nachweis $\not\Rightarrow$ instabil (anderes $V$ probieren).\quad
\textbullet\ $V$ muss pos.\ def.\ sein, sonst keine Aussage.\quad
\textbullet\ EZG ist nur der \emph{nachgewiesene} Teil (hängt von $V$ ab), nicht zwingend das gesamte.}

\bereichende
